\chapter{Embedded solids and cut cells}\label{ch:cutcells}
\chapterstatus[e805528147163a1c2fad4bbc5a5bf9d9106db836]{reference}
A fixed solid is embedded in the block-forest mesh without changing it: every
cell carries the volume fraction of the solid and, if it is cut, one
half-plane that approximates the solid inside it. This chapter states how the
solid is represented and reconstructed, how the open polygons, the one-sided
face traces, the shared apertures and the open staggered volumes are computed
from the half-planes, how the velocity is closed on faces that the solid
covers, and how the geometry is transferred when the mesh changes. The open
measures themselves and the operators built on them are defined in
\cref{sec:setting:open,sec:setting:divergence,sec:setting:elliptic}; phase
fractions in cut cells are those of \cref{sec:interface:fractions}. The
thesis introduces the open parts of cells and faces~\thesis{\S2.4} and refers
to the ComFLOW cut-cell method of Kleefsman et al.\ and Dröge and
Verstappen~\cite{verstappen2005} for the discretisation; the construction below
is the package's. The geometry is fixed between regrids, so time
discretisation is not applicable (\cref{sec:cutcells:time}).

\section{Continuous setting}\label{sec:cutcells:continuous}
The solid occupies the fixed region
\begin{equation}\label{eq:cutcells:domain}
  \domain^{\solid}=\set{\coord\in\domain}{\lset^{\solid}(\coord)\le0},\qquad
  \domain^{\open}=\domain\setminus\domain^{\solid},
\end{equation}
given by a level set $\lset^{\solid}$ (negative inside the solid; it need not
be a distance function). On the solid boundary the velocity satisfies the
no-penetration condition $\Velo\cdot\normal=0$; the tangential condition (no
slip or free slip) enters only the viscous stresses of \cref{ch:viscosity},
and the contact angle of an interface that meets the solid enters the
curvature (\cref{sec:interface:contact}). The divergence theorem on open cell
parts~\cref{eq:setting:cont-gauss} then involves only the open parts of the
cell faces.

\section{Solid reconstruction and open polygons}\label{sec:cutcells:solid}
The solid is an \apitype{InterfaceField} without geometry whose ``liquid'' is
the solid: its fractions $\solidfrac{c}$ are initialised from
$\lset^{\solid}$ by the level-set quadrature of \cref{sec:interface:init} and
reconstructed by \cref{sec:interface:plic,sec:interface:normals} (default
\code{:lhf}: LHF normals with the WY fallback, one-sided normals at domain
walls). Its fraction floor~\cref{eq:interface:floor} closes cells whose open
fraction is at most $\ffloor$ and removes solid slivers of at most $\ffloor$; no
other small-cell treatment is applied. In a cut cell, $0<\solidfrac{c}<1$, the
solid occupies the half-plane
\begin{equation}\label{eq:cutcells:solid}
  \localpart{\solid}{c}=\set{\lcoord\in\refcell}{\snormal{c}\cdot\lcoord\le\sconst{c}},
  \qquad \planearea(\snormal{c},\sconst{c};\refcell)=\solidfrac{c},
\end{equation}
so $\snormal{c}$ points into the fluid. The open polygon and the open fraction
are
\begin{equation}\label{eq:cutcells:open}
  \localpart{\open}{c}=\begin{cases}
    \refcell, & \solidfrac{c}=0,\\
    \set{\lcoord\in\refcell}{\snormal{c}\cdot\lcoord\ge\sconst{c}}, & 0<\solidfrac{c}<1,\\
    \emptyset, & \solidfrac{c}=1,
  \end{cases}
  \qquad
  \openfrac{c}=1-\solidfrac{c}=\meas{\localpart{\open}{c}},
\end{equation}
the last equality holding to the tolerance of the plane constant. Cells with
$\openfrac{c}>0$ are active; covered cells ($\openfrac{c}=0$) are inactive in
the operators of \cref{ch:discrete-setting}. The polygon is obtained by
clipping $\refcell$ with the half-plane (Sutherland--Hodgman~\cite{sutherland1974});
the PLIC segments of the solid and its triangulated reconstruction are
available for visualisation. Because each cut cell has its own half-plane, the
reconstructed boundary is in general discontinuous across faces
(\cref{fig:cutcells:apertures:traces}).

\emph{Check} (\code{figures/generate.jl}). On the cylinder fixture of
\cref{fig:cutcells:apertures} the area of $\localpart{\open}{c}$ equals
$\openfrac{c}$ to $\chkOpenFractionDiff$ in all \chkCutCells\ cut cells.

\begin{implnote}
\apitype{GeometryField}\code{(mesh, levelset; kwargs...)} (keywords of
\apitype{InterfaceField}) or \apitype{GeometryField}\code{(solid)}; fields
\code{solid}, \code{open_fraction} $=\openfrac{}$, \code{active_cells}:
\srcref{src/embedded/geometry.jl}{1}{82}; \apifun{refresh_geometry!},
\srcref{src/embedded/geometry.jl}{264}{318}; clipping
\srcref{src/embedded/geometry_utils.jl}{1}{94};
\apifun{plic_segments}, \apifun{plic_solid_geometry} and
\apitype{PLICSolidGeometry}, \srcref{src/visualization/plic_geometry.jl}{1}{73}.
Tests: \testref{test/indicator/geometry.jl}{499} (solid PLIC segments
converge with order above 1.49; triangulated area equals
$\sum_c\solidfrac{c}\meas{c}$), \testref{test/indicator/geometry.jl}{466}
(open fractions and apertures independent of the root partition),
\testref{test/two_phase/transport.jl}{93} (floor in cut cells).
\end{implnote}

\section{Face traces and shared apertures}\label{sec:cutcells:apertures}
Let $\localface{f}{c}$ be the image of a face segment $f\in\faceset{c}$ in the
local coordinates of $c$: $(\lcoord)_k=\pm\half$ on the side of $c$ that
contains $f$ ($k=\faxis{f}$, sign $+$ for $c=\cminus{f}$), and
$\rlo{t}\le(\lcoord)_t\le\rhi{t}$, $t\neq k$, the tangential extent of $f$ (a
half side at a level transition, \cref{eq:mesh:hanging}). The one-sided open
fraction of $f$ according to the half-plane of $c$ is
\begin{equation}\label{eq:cutcells:trace}
  \traceap{c}{f}=\frac{\meas{\localface{f}{c}\cap\localpart{\open}{c}}}{\meas{\localface{f}{c}}},
\end{equation}
one for an uncut and zero for a covered cell. On the segment the solid
condition $\snormal{c}\cdot\lcoord\le\sconst{c}$ is linear in $(\lcoord)_t$, so
its solid part is an interval, computed in closed form (the whole segment or
nothing if $\abs{(\snormal{c})_t}\le64\machineeps$). The shared aperture of
an interior face is the average of the two traces, zero if either cell is
covered~\cref{eq:setting:aperture}; a boundary face carries the trace of its
cell, $\aperture{f}=\traceap{\cbnd{f}}{f}$. One value per face is what makes
the open divergence conservative and the cut-cell gradient its negative
adjoint (\cref{sec:setting:divergence,sec:setting:elliptic}); the covered rule
closes faces on a grid-aligned solid boundary, where both cells are pure and
the average would otherwise be $\half$. The phase apertures of
\cref{sec:interface:apertures} split $\aperture{f}$ into liquid and gas parts.

\emph{Checks} (\code{figures/generate.jl}). On the cylinder fixture
(\chkCutFacesTotal\ interior faces, \chkCutFaces\ of them partially open), an
independent evaluation of~\cref{eq:cutcells:trace,eq:setting:aperture} from the
stored half-planes agrees with \code{geometry.aperture} to $\chkApertureDiff$
and with the boundary apertures to $\chkBoundaryApertureDiff$.

\begin{deviation}{Open apertures}
The thesis defines the open fraction and aperture as the sums of the phase
quantities, $\openfrac{}=\mean{\volfrac}$ and
$\aperture{}=\mean{a}$~\thesis{\S2.4}, with each phase aperture the average of
two one-sided traces~\thesis{eq.~(2.41)}. The code builds the open geometry
from the solid alone: $\openfrac{c}=1-\solidfrac{c}$, $\aperture{f}$ the
average of the two solid traces with the covered rule, and the phase apertures
split it, $\phaseaperture{\gas}{f}=\aperture{f}-\phaseaperture{\liq}{f}$.
On a face without a covered cell, $\aperture{f}$ equals the sum of the thesis' phase
apertures, $\half\sum_{c\in\cellset{f}}(\phasetrace{\liq}{c}{f}+\phasetrace{\gas}{c}{f})
=\half\sum_{c\in\cellset{f}}\traceap{c}{f}$ (the clamp
in~\cref{eq:interface:phase-aperture} is then inactive up to round-off). If a cell of
$\cellset{f}$ is covered, the thesis' sum is half of the open trace of the other cell
(one half at a grid-aligned solid boundary), the code's is zero.
\end{deviation}

\begin{implnote}
\code{_solid_line_fraction}, \code{_cell_solid_line_fraction},
\code{_face_local_tangent_interval}, \code{_cell_face_open_fraction}~\cref{eq:cutcells:trace},
\code{_interior_open_aperture}, \code{_boundary_open_aperture}:
\srcref{src/embedded/geometry.jl}{108}{182}; stored as \code{aperture} (order
of \code{mesh.faces}) and \code{boundary_aperture}. Tests:
\testref{test/indicator/geometry.jl}{427} (faces on a grid-aligned solid
boundary closed), \testref{test/indicator/geometry.jl}{1128} (liquid
aperture bounded by $\aperture{f}$ through two regrids).
\end{implnote}

\section{Open staggered volumes}\label{sec:cutcells:dual}
The open staggered volume~\cref{eq:setting:open-dual} bisects every open line
segment normal to the face~\cite{verstappen2005}. With
$\stripset{c}{f}$ the strip of $c$ over the full width of $c$ normal to $f$
and the tangential extent of $f$, whose local image is the rectangle
$[-\half,\half]\times[\rlo{t},\rhi{t}]$ (axes ordered $k,t$), each cell
contributes half of the open area of its strip,
\begin{equation}\label{eq:cutcells:dual}
  \opendualvol{f}=\half\sum_{c\in\cellset{f}}\meas{c}\,\roundpar{\rhi{t}-\rlo{t}}
    \roundpar{1-\planearea\roundpar{\snormal{c},\sconst{c};[-\half,\half]\times[\rlo{t},\rhi{t}]}},
\end{equation}
with $\planearea\equiv0$ for uncut and the contribution zero for covered
cells. A boundary face gets the half strip of its single cell; it is stored
but is not an unknown. Without a solid $\opendualvol{f}=\dualvol{f}$, but in
general $\opendualvol{f}\neq\aperture{f}\dualvol{f}$: for the face of
\cref{fig:cutcells:apertures:dual}, $\opendualvol{f}=\chkDualFace$ and
$\aperture{f}\dualvol{f}=\chkApertureDual$. The cut-cell elliptic operator
of \cref{sec:setting:elliptic} rescales its metric by
$(\dualvol{f}/\opendualvol{f})^{1/2}$.

\emph{Check} (\code{figures/generate.jl}). Evaluating~\cref{eq:cutcells:dual}
by clipping the strips independently reproduces \code{geometry.face_dual_volumes}
on all faces of the cylinder fixture to $\chkOpenDualDiff$.

% Figure: solid half-planes, face traces, shared apertures and an open
% staggered volume (Chapter 7). Polygons, intervals, apertures and volumes:
% generated/cutcell.tex (generate.jl cylinder fixture); labels hand-written.
\begin{figure}
  \centering
  \begin{subfigure}[b]{0.56\textwidth}
    \centering
    \begin{tikzpicture}[x=1.22cm,y=1.22cm]
      \def\fxccell#1#2#3#4{}\def\fxcsolid#1{}\def\fxcline#1#2#3#4{}%
      \def\fxctrace#1#2#3#4#5{}\def\fxcaperture#1#2#3#4{}\def\fxcstrip#1{}%
      \begin{scope}
        \clip (4,3) rectangle (8,7);
        % Pass 1: covered cells and the solid parts of cut cells.
        \def\fxccell#1#2#3#4{\ifnum#4=2 \fill[tr solid] (#1,#2) rectangle ++(1,1);\fi}%
        \def\fxcsolid#1{\fill[tr solid] #1;}%
        % Generated by docs/techref/figures/generate.jl; do not edit.
\gdef\fxCylX{4.3}\gdef\fxCylY{3.6}\gdef\fxCylR{2.6}
\fxccell{4}{3}{0.00}{2}
\fxccell{5}{3}{0.00}{2}
\fxccell{6}{3}{0.12}{1}
\fxcsolid{(6,3) -- (6.860509,3) -- (6.903338,4) -- (6,4) -- cycle}
\fxcline{6.860509}{3}{6.903338}{4}
\fxccell{7}{3}{1.00}{0}
\fxccell{4}{4}{0.00}{2}
\fxccell{5}{4}{0.00}{2}
\fxccell{6}{4}{0.28}{1}
\fxcsolid{(6,4) -- (7,4) -- (7,4.105177) -- (6.373381,5) -- (6,5) -- cycle}
\fxcline{7}{4.105177}{6.373381}{5}
\fxccell{7}{4}{1.00}{0}
\fxccell{4}{5}{0.00}{2}
\fxccell{5}{5}{0.13}{1}
\fxcsolid{(5,5) -- (6,5) -- (6,5.480591) -- (5.486078,6) -- (5,6) -- cycle}
\fxcline{6}{5.480591}{5.486078}{6}
\fxccell{6}{5}{0.85}{1}
\fxcsolid{(6,5.660692) -- (6,5) -- (6.462664,5) -- cycle}
\fxcline{6.462664}{5}{6}{5.660692}
\fxccell{7}{5}{1.00}{0}
\fxccell{4}{6}{0.82}{1}
\fxcsolid{(4,6.219014) -- (4,6) -- (5,6) -- (5,6.133053) -- cycle}
\fxcline{5}{6.133053}{4}{6.219014}
\fxccell{5}{6}{0.98}{1}
\fxcsolid{(5,6.098687) -- (5,6) -- (5.336959,6) -- cycle}
\fxcline{5.336959}{6}{5}{6.098687}
\fxccell{6}{6}{1.00}{0}
\fxccell{7}{6}{1.00}{0}
\fxctrace{1}{5}{5}{5}{-1}
\fxctrace{1}{5}{6}{6}{1}
\fxcaperture{1}{5}{5.5}{0.00}
\fxctrace{1}{5}{6.133053}{7}{-1}
\fxctrace{1}{5}{6.098687}{7}{1}
\fxcaperture{1}{5}{6.5}{0.88}
\fxctrace{1}{6}{3}{3}{-1}
\fxctrace{1}{6}{3}{3}{1}
\fxcaperture{1}{6}{3.5}{0.00}
\fxctrace{1}{6}{4}{4}{-1}
\fxctrace{1}{6}{5}{5}{1}
\fxcaperture{1}{6}{4.5}{0.00}
\fxctrace{1}{6}{5.480591}{6}{-1}
\fxctrace{1}{6}{5.660692}{6}{1}
\fxcaperture{1}{6}{5.5}{0.43}
\fxctrace{1}{6}{6}{7}{-1}
\fxctrace{1}{6}{6}{7}{1}
\fxcaperture{1}{6}{6.5}{1.00}
\fxctrace{1}{7}{3}{4}{-1}
\fxctrace{1}{7}{3}{4}{1}
\fxcaperture{1}{7}{3.5}{1.00}
\fxctrace{1}{7}{4.105177}{5}{-1}
\fxctrace{1}{7}{4}{5}{1}
\fxcaperture{1}{7}{4.5}{0.95}
\fxctrace{1}{7}{5}{6}{-1}
\fxctrace{1}{7}{5}{6}{1}
\fxcaperture{1}{7}{5.5}{1.00}
\fxctrace{2}{4}{6.903338}{7}{-1}
\fxctrace{2}{4}{7}{7}{1}
\fxcaperture{2}{6.5}{4}{0.05}
\fxctrace{2}{5}{5}{5}{-1}
\fxctrace{2}{5}{6}{6}{1}
\fxcaperture{2}{5.5}{5}{0.00}
\fxctrace{2}{5}{6.373381}{7}{-1}
\fxctrace{2}{5}{6.462664}{7}{1}
\fxcaperture{2}{6.5}{5}{0.58}
\fxctrace{2}{6}{4}{4}{-1}
\fxctrace{2}{6}{5}{5}{1}
\fxcaperture{2}{4.5}{6}{0.00}
\fxctrace{2}{6}{5.486078}{6}{-1}
\fxctrace{2}{6}{5.336959}{6}{1}
\fxcaperture{2}{5.5}{6}{0.59}
\fxctrace{2}{6}{6}{7}{-1}
\fxctrace{2}{6}{6}{7}{1}
\fxcaperture{2}{6.5}{6}{1.00}
\fxcstrip{(6,5.480591) -- (6,6) -- (5.486078,6) -- cycle}
\fxcstrip{(6.462664,5) -- (7,5) -- (7,6) -- (6,6) -- (6,5.660692) -- cycle}
\gdef\fxCdual{0.490}\gdef\fxCdualfull{1}\gdef\fxCaperture{0.43}\gdef\fxCfx{6}\gdef\fxCfy{5.5}

        % Pass 2: cells, solid lines and the exact boundary.
        \def\fxccell#1#2#3#4{\draw[tr cell] (#1,#2) rectangle ++(1,1);}%
        \def\fxcsolid#1{}%
        \def\fxcline#1#2#3#4{\draw[tr solid line] (#1,#2) -- (#3,#4);}%
        % Generated by docs/techref/figures/generate.jl; do not edit.
\gdef\fxCylX{4.3}\gdef\fxCylY{3.6}\gdef\fxCylR{2.6}
\fxccell{4}{3}{0.00}{2}
\fxccell{5}{3}{0.00}{2}
\fxccell{6}{3}{0.12}{1}
\fxcsolid{(6,3) -- (6.860509,3) -- (6.903338,4) -- (6,4) -- cycle}
\fxcline{6.860509}{3}{6.903338}{4}
\fxccell{7}{3}{1.00}{0}
\fxccell{4}{4}{0.00}{2}
\fxccell{5}{4}{0.00}{2}
\fxccell{6}{4}{0.28}{1}
\fxcsolid{(6,4) -- (7,4) -- (7,4.105177) -- (6.373381,5) -- (6,5) -- cycle}
\fxcline{7}{4.105177}{6.373381}{5}
\fxccell{7}{4}{1.00}{0}
\fxccell{4}{5}{0.00}{2}
\fxccell{5}{5}{0.13}{1}
\fxcsolid{(5,5) -- (6,5) -- (6,5.480591) -- (5.486078,6) -- (5,6) -- cycle}
\fxcline{6}{5.480591}{5.486078}{6}
\fxccell{6}{5}{0.85}{1}
\fxcsolid{(6,5.660692) -- (6,5) -- (6.462664,5) -- cycle}
\fxcline{6.462664}{5}{6}{5.660692}
\fxccell{7}{5}{1.00}{0}
\fxccell{4}{6}{0.82}{1}
\fxcsolid{(4,6.219014) -- (4,6) -- (5,6) -- (5,6.133053) -- cycle}
\fxcline{5}{6.133053}{4}{6.219014}
\fxccell{5}{6}{0.98}{1}
\fxcsolid{(5,6.098687) -- (5,6) -- (5.336959,6) -- cycle}
\fxcline{5.336959}{6}{5}{6.098687}
\fxccell{6}{6}{1.00}{0}
\fxccell{7}{6}{1.00}{0}
\fxctrace{1}{5}{5}{5}{-1}
\fxctrace{1}{5}{6}{6}{1}
\fxcaperture{1}{5}{5.5}{0.00}
\fxctrace{1}{5}{6.133053}{7}{-1}
\fxctrace{1}{5}{6.098687}{7}{1}
\fxcaperture{1}{5}{6.5}{0.88}
\fxctrace{1}{6}{3}{3}{-1}
\fxctrace{1}{6}{3}{3}{1}
\fxcaperture{1}{6}{3.5}{0.00}
\fxctrace{1}{6}{4}{4}{-1}
\fxctrace{1}{6}{5}{5}{1}
\fxcaperture{1}{6}{4.5}{0.00}
\fxctrace{1}{6}{5.480591}{6}{-1}
\fxctrace{1}{6}{5.660692}{6}{1}
\fxcaperture{1}{6}{5.5}{0.43}
\fxctrace{1}{6}{6}{7}{-1}
\fxctrace{1}{6}{6}{7}{1}
\fxcaperture{1}{6}{6.5}{1.00}
\fxctrace{1}{7}{3}{4}{-1}
\fxctrace{1}{7}{3}{4}{1}
\fxcaperture{1}{7}{3.5}{1.00}
\fxctrace{1}{7}{4.105177}{5}{-1}
\fxctrace{1}{7}{4}{5}{1}
\fxcaperture{1}{7}{4.5}{0.95}
\fxctrace{1}{7}{5}{6}{-1}
\fxctrace{1}{7}{5}{6}{1}
\fxcaperture{1}{7}{5.5}{1.00}
\fxctrace{2}{4}{6.903338}{7}{-1}
\fxctrace{2}{4}{7}{7}{1}
\fxcaperture{2}{6.5}{4}{0.05}
\fxctrace{2}{5}{5}{5}{-1}
\fxctrace{2}{5}{6}{6}{1}
\fxcaperture{2}{5.5}{5}{0.00}
\fxctrace{2}{5}{6.373381}{7}{-1}
\fxctrace{2}{5}{6.462664}{7}{1}
\fxcaperture{2}{6.5}{5}{0.58}
\fxctrace{2}{6}{4}{4}{-1}
\fxctrace{2}{6}{5}{5}{1}
\fxcaperture{2}{4.5}{6}{0.00}
\fxctrace{2}{6}{5.486078}{6}{-1}
\fxctrace{2}{6}{5.336959}{6}{1}
\fxcaperture{2}{5.5}{6}{0.59}
\fxctrace{2}{6}{6}{7}{-1}
\fxctrace{2}{6}{6}{7}{1}
\fxcaperture{2}{6.5}{6}{1.00}
\fxcstrip{(6,5.480591) -- (6,6) -- (5.486078,6) -- cycle}
\fxcstrip{(6.462664,5) -- (7,5) -- (7,6) -- (6,6) -- (6,5.660692) -- cycle}
\gdef\fxCdual{0.490}\gdef\fxCdualfull{1}\gdef\fxCaperture{0.43}\gdef\fxCfx{6}\gdef\fxCfy{5.5}

        \draw[tr exact] (\fxCylX,\fxCylY) circle (\fxCylR);
      \end{scope}
      % Pass 3: open intervals of each face seen from either cell (offset to
      % the side of that cell), shared apertures, open fractions.
      \def\fxccell#1#2#3#4{\ifnum#4=1 \node[tr value,font=\scriptsize] at (#1+0.5,#2+0.5) {#3};\fi}%
      \def\fxcline#1#2#3#4{}%
      \def\fxctrace#1#2#3#4#5{%
        \ifdim#3pt<#4pt
          \ifnum#1=1 \draw[tr trace,line width=1.1pt] (#2+0.07*#5,#3) -- (#2+0.07*#5,#4);
          \else      \draw[tr trace,line width=1.1pt] (#3,#2+0.07*#5) -- (#4,#2+0.07*#5);\fi
        \fi}%
      \def\fxcaperture#1#2#3#4{%
        \ifdim#4pt>0pt\ifdim#4pt<1pt
          \ifnum#1=1 \node[tr tiny,anchor=west,inner sep=0.6pt,text=trgreen!40!black] at (#2+0.1,#3+0.32) {#4};
          \else      \node[tr tiny,anchor=south,inner sep=0.6pt,text=trgreen!40!black] at (#2+0.3,#3+0.08) {#4};\fi
        \fi\fi}%
      % Generated by docs/techref/figures/generate.jl; do not edit.
\gdef\fxCylX{4.3}\gdef\fxCylY{3.6}\gdef\fxCylR{2.6}
\fxccell{4}{3}{0.00}{2}
\fxccell{5}{3}{0.00}{2}
\fxccell{6}{3}{0.12}{1}
\fxcsolid{(6,3) -- (6.860509,3) -- (6.903338,4) -- (6,4) -- cycle}
\fxcline{6.860509}{3}{6.903338}{4}
\fxccell{7}{3}{1.00}{0}
\fxccell{4}{4}{0.00}{2}
\fxccell{5}{4}{0.00}{2}
\fxccell{6}{4}{0.28}{1}
\fxcsolid{(6,4) -- (7,4) -- (7,4.105177) -- (6.373381,5) -- (6,5) -- cycle}
\fxcline{7}{4.105177}{6.373381}{5}
\fxccell{7}{4}{1.00}{0}
\fxccell{4}{5}{0.00}{2}
\fxccell{5}{5}{0.13}{1}
\fxcsolid{(5,5) -- (6,5) -- (6,5.480591) -- (5.486078,6) -- (5,6) -- cycle}
\fxcline{6}{5.480591}{5.486078}{6}
\fxccell{6}{5}{0.85}{1}
\fxcsolid{(6,5.660692) -- (6,5) -- (6.462664,5) -- cycle}
\fxcline{6.462664}{5}{6}{5.660692}
\fxccell{7}{5}{1.00}{0}
\fxccell{4}{6}{0.82}{1}
\fxcsolid{(4,6.219014) -- (4,6) -- (5,6) -- (5,6.133053) -- cycle}
\fxcline{5}{6.133053}{4}{6.219014}
\fxccell{5}{6}{0.98}{1}
\fxcsolid{(5,6.098687) -- (5,6) -- (5.336959,6) -- cycle}
\fxcline{5.336959}{6}{5}{6.098687}
\fxccell{6}{6}{1.00}{0}
\fxccell{7}{6}{1.00}{0}
\fxctrace{1}{5}{5}{5}{-1}
\fxctrace{1}{5}{6}{6}{1}
\fxcaperture{1}{5}{5.5}{0.00}
\fxctrace{1}{5}{6.133053}{7}{-1}
\fxctrace{1}{5}{6.098687}{7}{1}
\fxcaperture{1}{5}{6.5}{0.88}
\fxctrace{1}{6}{3}{3}{-1}
\fxctrace{1}{6}{3}{3}{1}
\fxcaperture{1}{6}{3.5}{0.00}
\fxctrace{1}{6}{4}{4}{-1}
\fxctrace{1}{6}{5}{5}{1}
\fxcaperture{1}{6}{4.5}{0.00}
\fxctrace{1}{6}{5.480591}{6}{-1}
\fxctrace{1}{6}{5.660692}{6}{1}
\fxcaperture{1}{6}{5.5}{0.43}
\fxctrace{1}{6}{6}{7}{-1}
\fxctrace{1}{6}{6}{7}{1}
\fxcaperture{1}{6}{6.5}{1.00}
\fxctrace{1}{7}{3}{4}{-1}
\fxctrace{1}{7}{3}{4}{1}
\fxcaperture{1}{7}{3.5}{1.00}
\fxctrace{1}{7}{4.105177}{5}{-1}
\fxctrace{1}{7}{4}{5}{1}
\fxcaperture{1}{7}{4.5}{0.95}
\fxctrace{1}{7}{5}{6}{-1}
\fxctrace{1}{7}{5}{6}{1}
\fxcaperture{1}{7}{5.5}{1.00}
\fxctrace{2}{4}{6.903338}{7}{-1}
\fxctrace{2}{4}{7}{7}{1}
\fxcaperture{2}{6.5}{4}{0.05}
\fxctrace{2}{5}{5}{5}{-1}
\fxctrace{2}{5}{6}{6}{1}
\fxcaperture{2}{5.5}{5}{0.00}
\fxctrace{2}{5}{6.373381}{7}{-1}
\fxctrace{2}{5}{6.462664}{7}{1}
\fxcaperture{2}{6.5}{5}{0.58}
\fxctrace{2}{6}{4}{4}{-1}
\fxctrace{2}{6}{5}{5}{1}
\fxcaperture{2}{4.5}{6}{0.00}
\fxctrace{2}{6}{5.486078}{6}{-1}
\fxctrace{2}{6}{5.336959}{6}{1}
\fxcaperture{2}{5.5}{6}{0.59}
\fxctrace{2}{6}{6}{7}{-1}
\fxctrace{2}{6}{6}{7}{1}
\fxcaperture{2}{6.5}{6}{1.00}
\fxcstrip{(6,5.480591) -- (6,6) -- (5.486078,6) -- cycle}
\fxcstrip{(6.462664,5) -- (7,5) -- (7,6) -- (6,6) -- (6,5.660692) -- cycle}
\gdef\fxCdual{0.490}\gdef\fxCdualfull{1}\gdef\fxCaperture{0.43}\gdef\fxCfx{6}\gdef\fxCfy{5.5}

      \draw[line width=0.7pt] (\fxCfx,\fxCfy-0.5) -- (\fxCfx,\fxCfy+0.5);
      \node[tr small,anchor=south] at (\fxCfx,\fxCfy+0.5) {$f$};
    \end{tikzpicture}
    \caption{Solid half-planes (hatched) and open fractions $\openfrac{c}$;
      green: open part of each face seen from either side, number: $\aperture{f}$.}
    \label{fig:cutcells:apertures:traces}
  \end{subfigure}%
  \hfill
  \begin{subfigure}[b]{0.41\textwidth}
    \centering
    \begin{tikzpicture}[x=2.2cm,y=2.2cm]
      \def\fxccell#1#2#3#4{}\def\fxcsolid#1{}\def\fxcline#1#2#3#4{}%
      \def\fxctrace#1#2#3#4#5{}\def\fxcaperture#1#2#3#4{}\def\fxcstrip#1{}%
      \begin{scope}
        \clip (\fxCfx-1,\fxCfy-0.5) rectangle (\fxCfx+1,\fxCfy+0.5);
        \def\fxcsolid#1{\fill[tr solid] #1;}%
        \def\fxcstrip#1{\fill[tr dual fill] #1;}%
        % Generated by docs/techref/figures/generate.jl; do not edit.
\gdef\fxCylX{4.3}\gdef\fxCylY{3.6}\gdef\fxCylR{2.6}
\fxccell{4}{3}{0.00}{2}
\fxccell{5}{3}{0.00}{2}
\fxccell{6}{3}{0.12}{1}
\fxcsolid{(6,3) -- (6.860509,3) -- (6.903338,4) -- (6,4) -- cycle}
\fxcline{6.860509}{3}{6.903338}{4}
\fxccell{7}{3}{1.00}{0}
\fxccell{4}{4}{0.00}{2}
\fxccell{5}{4}{0.00}{2}
\fxccell{6}{4}{0.28}{1}
\fxcsolid{(6,4) -- (7,4) -- (7,4.105177) -- (6.373381,5) -- (6,5) -- cycle}
\fxcline{7}{4.105177}{6.373381}{5}
\fxccell{7}{4}{1.00}{0}
\fxccell{4}{5}{0.00}{2}
\fxccell{5}{5}{0.13}{1}
\fxcsolid{(5,5) -- (6,5) -- (6,5.480591) -- (5.486078,6) -- (5,6) -- cycle}
\fxcline{6}{5.480591}{5.486078}{6}
\fxccell{6}{5}{0.85}{1}
\fxcsolid{(6,5.660692) -- (6,5) -- (6.462664,5) -- cycle}
\fxcline{6.462664}{5}{6}{5.660692}
\fxccell{7}{5}{1.00}{0}
\fxccell{4}{6}{0.82}{1}
\fxcsolid{(4,6.219014) -- (4,6) -- (5,6) -- (5,6.133053) -- cycle}
\fxcline{5}{6.133053}{4}{6.219014}
\fxccell{5}{6}{0.98}{1}
\fxcsolid{(5,6.098687) -- (5,6) -- (5.336959,6) -- cycle}
\fxcline{5.336959}{6}{5}{6.098687}
\fxccell{6}{6}{1.00}{0}
\fxccell{7}{6}{1.00}{0}
\fxctrace{1}{5}{5}{5}{-1}
\fxctrace{1}{5}{6}{6}{1}
\fxcaperture{1}{5}{5.5}{0.00}
\fxctrace{1}{5}{6.133053}{7}{-1}
\fxctrace{1}{5}{6.098687}{7}{1}
\fxcaperture{1}{5}{6.5}{0.88}
\fxctrace{1}{6}{3}{3}{-1}
\fxctrace{1}{6}{3}{3}{1}
\fxcaperture{1}{6}{3.5}{0.00}
\fxctrace{1}{6}{4}{4}{-1}
\fxctrace{1}{6}{5}{5}{1}
\fxcaperture{1}{6}{4.5}{0.00}
\fxctrace{1}{6}{5.480591}{6}{-1}
\fxctrace{1}{6}{5.660692}{6}{1}
\fxcaperture{1}{6}{5.5}{0.43}
\fxctrace{1}{6}{6}{7}{-1}
\fxctrace{1}{6}{6}{7}{1}
\fxcaperture{1}{6}{6.5}{1.00}
\fxctrace{1}{7}{3}{4}{-1}
\fxctrace{1}{7}{3}{4}{1}
\fxcaperture{1}{7}{3.5}{1.00}
\fxctrace{1}{7}{4.105177}{5}{-1}
\fxctrace{1}{7}{4}{5}{1}
\fxcaperture{1}{7}{4.5}{0.95}
\fxctrace{1}{7}{5}{6}{-1}
\fxctrace{1}{7}{5}{6}{1}
\fxcaperture{1}{7}{5.5}{1.00}
\fxctrace{2}{4}{6.903338}{7}{-1}
\fxctrace{2}{4}{7}{7}{1}
\fxcaperture{2}{6.5}{4}{0.05}
\fxctrace{2}{5}{5}{5}{-1}
\fxctrace{2}{5}{6}{6}{1}
\fxcaperture{2}{5.5}{5}{0.00}
\fxctrace{2}{5}{6.373381}{7}{-1}
\fxctrace{2}{5}{6.462664}{7}{1}
\fxcaperture{2}{6.5}{5}{0.58}
\fxctrace{2}{6}{4}{4}{-1}
\fxctrace{2}{6}{5}{5}{1}
\fxcaperture{2}{4.5}{6}{0.00}
\fxctrace{2}{6}{5.486078}{6}{-1}
\fxctrace{2}{6}{5.336959}{6}{1}
\fxcaperture{2}{5.5}{6}{0.59}
\fxctrace{2}{6}{6}{7}{-1}
\fxctrace{2}{6}{6}{7}{1}
\fxcaperture{2}{6.5}{6}{1.00}
\fxcstrip{(6,5.480591) -- (6,6) -- (5.486078,6) -- cycle}
\fxcstrip{(6.462664,5) -- (7,5) -- (7,6) -- (6,6) -- (6,5.660692) -- cycle}
\gdef\fxCdual{0.490}\gdef\fxCdualfull{1}\gdef\fxCaperture{0.43}\gdef\fxCfx{6}\gdef\fxCfy{5.5}

        \def\fxcsolid#1{}\def\fxcstrip#1{}%
        \def\fxcline#1#2#3#4{\draw[tr solid line] (#1,#2) -- (#3,#4);}%
        % Generated by docs/techref/figures/generate.jl; do not edit.
\gdef\fxCylX{4.3}\gdef\fxCylY{3.6}\gdef\fxCylR{2.6}
\fxccell{4}{3}{0.00}{2}
\fxccell{5}{3}{0.00}{2}
\fxccell{6}{3}{0.12}{1}
\fxcsolid{(6,3) -- (6.860509,3) -- (6.903338,4) -- (6,4) -- cycle}
\fxcline{6.860509}{3}{6.903338}{4}
\fxccell{7}{3}{1.00}{0}
\fxccell{4}{4}{0.00}{2}
\fxccell{5}{4}{0.00}{2}
\fxccell{6}{4}{0.28}{1}
\fxcsolid{(6,4) -- (7,4) -- (7,4.105177) -- (6.373381,5) -- (6,5) -- cycle}
\fxcline{7}{4.105177}{6.373381}{5}
\fxccell{7}{4}{1.00}{0}
\fxccell{4}{5}{0.00}{2}
\fxccell{5}{5}{0.13}{1}
\fxcsolid{(5,5) -- (6,5) -- (6,5.480591) -- (5.486078,6) -- (5,6) -- cycle}
\fxcline{6}{5.480591}{5.486078}{6}
\fxccell{6}{5}{0.85}{1}
\fxcsolid{(6,5.660692) -- (6,5) -- (6.462664,5) -- cycle}
\fxcline{6.462664}{5}{6}{5.660692}
\fxccell{7}{5}{1.00}{0}
\fxccell{4}{6}{0.82}{1}
\fxcsolid{(4,6.219014) -- (4,6) -- (5,6) -- (5,6.133053) -- cycle}
\fxcline{5}{6.133053}{4}{6.219014}
\fxccell{5}{6}{0.98}{1}
\fxcsolid{(5,6.098687) -- (5,6) -- (5.336959,6) -- cycle}
\fxcline{5.336959}{6}{5}{6.098687}
\fxccell{6}{6}{1.00}{0}
\fxccell{7}{6}{1.00}{0}
\fxctrace{1}{5}{5}{5}{-1}
\fxctrace{1}{5}{6}{6}{1}
\fxcaperture{1}{5}{5.5}{0.00}
\fxctrace{1}{5}{6.133053}{7}{-1}
\fxctrace{1}{5}{6.098687}{7}{1}
\fxcaperture{1}{5}{6.5}{0.88}
\fxctrace{1}{6}{3}{3}{-1}
\fxctrace{1}{6}{3}{3}{1}
\fxcaperture{1}{6}{3.5}{0.00}
\fxctrace{1}{6}{4}{4}{-1}
\fxctrace{1}{6}{5}{5}{1}
\fxcaperture{1}{6}{4.5}{0.00}
\fxctrace{1}{6}{5.480591}{6}{-1}
\fxctrace{1}{6}{5.660692}{6}{1}
\fxcaperture{1}{6}{5.5}{0.43}
\fxctrace{1}{6}{6}{7}{-1}
\fxctrace{1}{6}{6}{7}{1}
\fxcaperture{1}{6}{6.5}{1.00}
\fxctrace{1}{7}{3}{4}{-1}
\fxctrace{1}{7}{3}{4}{1}
\fxcaperture{1}{7}{3.5}{1.00}
\fxctrace{1}{7}{4.105177}{5}{-1}
\fxctrace{1}{7}{4}{5}{1}
\fxcaperture{1}{7}{4.5}{0.95}
\fxctrace{1}{7}{5}{6}{-1}
\fxctrace{1}{7}{5}{6}{1}
\fxcaperture{1}{7}{5.5}{1.00}
\fxctrace{2}{4}{6.903338}{7}{-1}
\fxctrace{2}{4}{7}{7}{1}
\fxcaperture{2}{6.5}{4}{0.05}
\fxctrace{2}{5}{5}{5}{-1}
\fxctrace{2}{5}{6}{6}{1}
\fxcaperture{2}{5.5}{5}{0.00}
\fxctrace{2}{5}{6.373381}{7}{-1}
\fxctrace{2}{5}{6.462664}{7}{1}
\fxcaperture{2}{6.5}{5}{0.58}
\fxctrace{2}{6}{4}{4}{-1}
\fxctrace{2}{6}{5}{5}{1}
\fxcaperture{2}{4.5}{6}{0.00}
\fxctrace{2}{6}{5.486078}{6}{-1}
\fxctrace{2}{6}{5.336959}{6}{1}
\fxcaperture{2}{5.5}{6}{0.59}
\fxctrace{2}{6}{6}{7}{-1}
\fxctrace{2}{6}{6}{7}{1}
\fxcaperture{2}{6.5}{6}{1.00}
\fxcstrip{(6,5.480591) -- (6,6) -- (5.486078,6) -- cycle}
\fxcstrip{(6.462664,5) -- (7,5) -- (7,6) -- (6,6) -- (6,5.660692) -- cycle}
\gdef\fxCdual{0.490}\gdef\fxCdualfull{1}\gdef\fxCaperture{0.43}\gdef\fxCfx{6}\gdef\fxCfy{5.5}

      \end{scope}
      \draw[tr cell] (\fxCfx-1,\fxCfy-0.5) rectangle (\fxCfx+1,\fxCfy+0.5);
      \draw[tr dual] (\fxCfx-0.5,\fxCfy-0.5) rectangle (\fxCfx+0.5,\fxCfy+0.5);
      \draw[tr face] (\fxCfx,\fxCfy-0.5) -- (\fxCfx,\fxCfy+0.5);
      \node[tr small,anchor=south] at (\fxCfx,\fxCfy+0.5) {$f$};
      \node[tr small,anchor=north] at (\fxCfx-0.5,\fxCfy-0.5) {$\cminus{f}$};
      \node[tr small,anchor=north] at (\fxCfx+0.5,\fxCfy-0.5) {$\cplus{f}$};
      \node[tr small,align=center,anchor=north] at (\fxCfx,\fxCfy-0.72)
        {$\opendualvol{f}=\chkDualFace$, \quad $\aperture{f}\dualvol{f}=\chkApertureDual$};
    \end{tikzpicture}
    \caption{Open staggered volume of $f$: half of the open parts (green) of
      the two strips $\stripset{c}{f}$; dashed: $\dualcv{f}$.}
    \label{fig:cutcells:apertures:dual}
  \end{subfigure}
  \caption{Cut-cell geometry on the cylinder fixture of \code{figures/generate.jl}
    (solid disc of radius 2.6 in $8\times8$ unit cells; window
    $[4,8]\times[3,7]$; dashed red: the exact boundary). Each cut
    cell carries its own solid half-plane~\cref{eq:cutcells:solid}, so the two
    traces~\cref{eq:cutcells:trace} of a face differ; the shared
    aperture~\cref{eq:setting:aperture} averages them and is zero next to a
    covered cell. (b) The face $f$ marked in (a): the open staggered
    volume~\cref{eq:setting:open-dual} is not $\aperture{f}\dualvol{f}$.}
  \label{fig:cutcells:apertures}
\end{figure}


\begin{implnote}
\code{_open_subcell_area}, \code{_geometry_dual_volumes} and
\code{_geometry_boundary_dual_volumes}:
\srcref{src/embedded/geometry.jl}{184}{262}; stored per axis in
\code{dual_volumes} (operator order), \code{face_dual_volumes} (order of
\code{mesh.faces}) and \code{boundary_dual_volumes}.
\end{implnote}

\section{Velocity closure and operators on cut cells}\label{sec:cutcells:operators}
Faces that the solid closes carry the stationary-wall value,
\begin{equation}\label{eq:cutcells:closed}
  \velo_f\leftarrow0\quad\text{for every }f\in\faces\text{ with }\aperture{f}=0,
\end{equation}
applied wherever the time step enforces the velocity boundary conditions
(\cref{ch:overview}) and before and after the cut-cell projection (for the two
phase velocities, values marked absent are kept). On a partially open face $\velo_f$
is the velocity of the open part: the solid part of $\partial\opencell{c}$
carries no flux in the open divergence~\cref{eq:setting:open-divergence},
which is the discrete form of $\Velo\cdot\normal=0$ there. The cut-cell
operators use only the quantities of this chapter: the open divergence and
pressure gradient (\cref{sec:setting:divergence,sec:setting:elliptic}), the
upwind transport of a cell scalar with fluxes
$\dt\,\meas{f}\aperture{f}\velo_f$ and updates divided by $\meas{\opencell{c}}$
(\cref{ch:momentum-transport}), the geometric VOF fluxes of
\cref{ch:vof-transport} and the wall stresses of \cref{ch:viscosity}. Open
cells remain active however small they are; there is no merging or flux
redistribution; the geometric transport bounds its time step with the open
volumes (\cref{ch:vof-transport,ch:stability-accuracy}).

\begin{implnote}
\apifun{enforce_geometry_no_slip!}\code{(velocity, geometry; preserve_absent=false)}~\cref{eq:cutcells:closed},
\apifun{divergence!} and \apifun{advect!} with a \apitype{GeometryField}:
\srcref{src/embedded/operators.jl}{1}{119}; calls in the time step
\srcref{src/single_phase/stepping.jl}{3}{12},
\srcref{src/two_phase/two_velocity/stepping.jl}{24}{28} and in the cut-cell
projection \srcref{src/embedded/operators.jl}{316}{376}; open-volume CFL
bound of the transport \srcref{src/vof/workspace.jl}{282}{297}.
\end{implnote}

\section{Regrid and rebuild}\label{sec:cutcells:time}
Time discretisation is not applicable: the geometry depends on the mesh only.
When the leaf set changes, the solid fractions are transferred geometrically.
A child $c_r$, $r=1,\dots,4$, of a refined cut cell $c$ receives the solid
area of the parent's half-plane over its quadrant $\refcell_r\subset\refcell$
of the parent's local coordinates,
\begin{equation}\label{eq:cutcells:prolong}
  \solidfrac{c_r}=4\,\meas{\refcell_r\cap\localpart{\solid}{c}},\qquad
  \solidfrac{c}=\tfrac14\sum_{r=1}^{4}\solidfrac{c_r}\ \ \text{(coarsening)},
\end{equation}
which conserves the solid volume exactly up to rounding. Then the solid is
reconstructed and every quantity of
\cref{sec:cutcells:solid,sec:cutcells:apertures,sec:cutcells:dual} is
recomputed; a revision counter makes the phase interface reconstruct its
cut cells again. The level set $\lset^{\solid}$ is not evaluated again, so
refining near the solid does not make its representation more accurate than
the half-planes from which the children inherit their fractions; construct
the \apitype{GeometryField} on a mesh that already resolves the solid.

\begin{deviation}{ComFLOW boundary features}
The thesis' ComFLOW-specific free-surface and wave-generation boundary
conditions are not implemented.
\end{deviation}

\begin{implnote}
\code{GeometryProlongation} (the interface prolongation applied to the solid,
then \apifun{refresh_geometry!}) and \apifun{register_adaptive_field!}:
\srcref{src/embedded/geometry.jl}{320}{363}; the child clipping
\srcref{src/interface/interface_field.jl}{571}{629}. Tests:
\testref{test/indicator/geometry.jl}{1128} (solid and phase volumes conserved
to $5\times10^{-14}$ through refinement and coarsening),
\testref{test/indicator/geometry.jl}{773} (operators rebuilt after a regrid).
\end{implnote}
